% Section 6: Discussion
\section{Discussion}
\label{sec:discussion}

\subsection{Summary}

We have shown that the incomplete-data structure of a parallel system
observation is component-level left-censoring with latent identity,
not masked cause of failure.  This censoring equivalence
(\Cref{prop:censoring-equiv}) provides the conceptual foundation for
a closed-form MLE via inclusion-exclusion, an EM algorithm with
analytical E- and M-steps, and an elementary identifiability proof for
exponential components.  The simulations of \Cref{sec:simulations}
confirm that the MLE is consistent with nominal Wald coverage, and
reveal a striking information asymmetry: the precision of
$\hat\lambda_j$ is governed by $P(T_j = T)$, the probability that
component~$j$ determines the system lifetime.

\subsection{The information asymmetry}

The quantity $P(T_j = T)$ emerges as the organizing principle for
parallel system estimation.  Slow components (small~$\lambda_j$) fail
last with high probability, so their failure times are observed nearly
exactly in most system observations.  Fast components (large~$\lambda_j$)
fail early and are deeply left-censored by the system structure, yielding
RMSE ratios exceeding~$150{:}1$ in the five-component case.

This has a direct practical implication.  If reliable estimates of all
component rates are needed, system-level observation alone is
insufficient for fast components.  Additional component-level
information---through periodic inspection (Scheme~1), partial
monitoring (Scheme~3), or degradation signatures---is required.  The
Fisher information analysis of \Cref{sec:sim-fisher} quantifies
this: the per-component efficiency under Scheme~0 is approximately
proportional to $P(T_j = T)$, and the joint efficiency (determinant
ratio) collapses exponentially with~$m$, from $5.8 \times 10^{-2}$ at
$m = 2$ to $6.1 \times 10^{-10}$ at $m = 5$.  Extending this
comparison to Scheme~1 (periodic inspection) and other intermediate
observation schemes is a natural direction for guiding monitoring
strategy.

The asymmetry mirrors the dual situation in series systems.  There,
the weakest-link component ($\argmin_j T_j$) dominates the system
lifetime and is well-estimated from system data, while stronger
components are right-censored.  In parallel systems, the
strongest-link component ($\argmax_j T_j$) dominates, and weaker
components are left-censored.  The censoring direction reverses, but
the structural pattern is the same: the component that determines the
system lifetime is the one about which the most is learned.

\subsection{Limitations}

\paragraph{Exponential assumption.}
The constant-hazard assumption of the exponential distribution is
restrictive for many applications.  The inclusion-exclusion expansion
and the closed-form integrals (\Cref{lem:ie,lem:integral-wj}) depend
on the exponential form of the component CDF.  Extension to Weibull
or gamma components would preserve the general structure of the
density decomposition (\Cref{lem:density-decomp}) and the EM
framework (\Cref{prop:em-algorithm}), but the closed-form integrals
would be replaced by numerical quadrature, and the identifiability
proof would require different tail analysis.

\paragraph{Computational scaling.}
The inclusion-exclusion expansion has $2^m$ terms, limiting practical
applicability to systems with roughly $m \leq 15$--$20$ components.
For larger systems, Monte Carlo estimation of the density, recursive
decomposition of the inclusion-exclusion sum, or restriction to i.i.d.\
component models would be needed.

\paragraph{Independence.}
The independence assumption underlies both the product-form CDF and the
identifiability proof (\Cref{thm:identifiability}).  This is not merely
a convenience: the \citet{BasuGhosh1980} identifiability result
\emph{requires} independence.  Under general dependence, the
distribution of $\max(T_1, \ldots, T_m)$ no longer uniquely determines
the marginal component distributions---the same system CDF can arise
from different combinations of marginals and copula.  This is the
parallel-system analogue of the \citet{Tsiatis1975} non-identifiability
for series systems, but it appears to be unstated in the literature:
recent work on identifiability under dependence
\citep{HiabuLoWilke2025} addresses the series case via exclusion
restrictions, while the parallel case under dependence remains open.
Copula extensions \citep{NavarroNgBalakrishnan2012} could model
dependence but would require either structural assumptions on the
copula family or additional observation beyond Scheme~0.

\paragraph{Interval-censored optimization.}
The simulation study (\Cref{sec:sim-censoring}) revealed convergence
difficulties for purely interval-censored system data under direct
maximization.  The EM algorithm may provide better numerical behavior
in this setting, but a systematic comparison was beyond the scope of
the present study.

\subsection{Extension to $k$-out-of-$n$ systems}
\label{sec:k-out-of-n}

The parallel system ($1$-out-of-$n$) is the extreme case of a
$k$-out-of-$n$ system: the system functions as long as at least $k$ of
its $n$ components survive.  Equivalently, the system fails when the
$(n - k + 1)$-th component fails:
\begin{equation}\label{eq:kofn-lifetime}
  T = T_{(n-k+1)},
\end{equation}
where $T_{(1)} \leq T_{(2)} \leq \cdots \leq T_{(n)}$ denote the order
statistics of the component lifetimes.  The parallel system has $k = 1$
(system lifetime $T = T_{(n)}$, the maximum), and the series system has
$k = n$ (system lifetime $T = T_{(1)}$, the minimum).

The censoring equivalence (\Cref{prop:censoring-equiv}) generalizes
naturally to this setting.

\begin{proposition}[$k$-out-of-$n$ censoring equivalence]
\label{prop:kofn-censoring}
Consider independent, continuous component lifetimes
$T_1, \ldots, T_n$ and a $k$-out-of-$n$ system with failure
time $T = T_{(n-k+1)}$.  Observing the system lifetime~$T$ under
Scheme~0 is informationally equivalent to the following component-level
censoring scheme:
\begin{enumerate}
  \item[\textup{(a)}] One component---the $(n-k+1)$-th to fail---is
    observed exactly: $T_J = T$ for a latent index~$J$.
  \item[\textup{(b)}] The $n - k$ components that failed before the system
    are left-censored at~$T$: for these indices~$j$, we know only that
    $T_j < T$.
  \item[\textup{(c)}] The $k - 1$ components still functioning at system
    failure are right-censored at~$T$: for these indices~$j$, we know
    only that $T_j > T$.
  \item[\textup{(d)}] The identities of the components in each group are
    latent.
\end{enumerate}
\end{proposition}

\begin{proof}
When the system fails at time $T = T_{(n-k+1)} = t$, exactly $n - k + 1$
components have failed (those with $T_j \leq t$) and exactly $k - 1$
components survive (those with $T_j > t$).  Among the failed components,
exactly one satisfies $T_j = t$ almost surely (by continuity of the
component distributions), and the remaining $n - k$ have $T_j < t$.
Under Scheme~0, the system-level observation~$T = t$ reveals only this
count structure, not which components belong to which group.  This yields
one exact, $n - k$ left-censored, and $k - 1$ right-censored component
observations with latent assignment.
\end{proof}

\Cref{prop:kofn-censoring} recovers \Cref{prop:censoring-equiv} when
$k = 1$: the parallel system has $n - 1$ left-censored and zero
right-censored observations.  At the other extreme, $k = n$ gives the
series system: zero left-censored and $n - 1$ right-censored
observations---exactly the competing risks structure underlying the
classical masked cause-of-failure problem.

\begin{remark}[Information hierarchy]
\label{rmk:info-hierarchy}
The proposition reveals a monotone information structure as $k$
varies.  Decreasing~$k$ (from series toward parallel) converts
right-censored components into left-censored ones, but the total number
of censored components remains $n - 1$ regardless of~$k$.  What changes
is the \emph{direction}: left-censoring (parallel) and right-censoring
(series) carry different amounts of information, giving rise to
the identifiability contrast between the two extremes
(\Cref{cor:duality}).  For intermediate~$k$, both censoring
directions coexist, and the information content interpolates.
\end{remark}

For the discrete analogue (geometric component lifetimes),
\citet{Jasinski2023} and \citet{Jasinski2024} have derived closed-form
MLEs for $k$-out-of-$n$ systems with homogeneous and heterogeneous
components, respectively.  \Cref{prop:kofn-censoring} provides the
continuous counterpart, establishing the censoring structure that
underlies estimation for any parametric family.

\subsection{Related developments}

The censoring framework developed here complements several recent
research directions.  \citet{Cramer2023} provides a comprehensive
review of ordered and censored lifetime data in reliability, situating
progressive censoring models (including the coherent system
representation of \citealt{CramerNavarro2015}) within a unified
framework.  \citet{QuNgMoon2024a,QuNgMoon2024b} develop nonparametric
tests for component distribution homogeneity from system data with known
signatures---a testing complement to our estimation approach.
\citet{BiswasGangulyMitra2025} propose semiparametric (piecewise-linear
cumulative hazard) models for load-sharing parallel systems, providing a
flexible alternative to the parametric exponential model studied here.
These directions---nonparametric testing, semiparametric models, and the
parametric censoring-based MLE developed in the present paper---represent
complementary approaches to the same fundamental problem of learning
about components from system-level data.

A structural connection also exists with the one-shot device testing
literature \citep{BalakrishnanSo2025}, where each test unit is
destroyed upon inspection and only pass/fail status at the inspection
time is recorded.  For a parallel system observed at a single time
point (current-status data), the Scheme~0 observation reduces to a
binary outcome---structurally identical to the one-shot device model
with multiple competing failure modes.  The two literatures have
developed largely independently; the censoring equivalence developed
here provides a bridge.

\subsection{Future work}

Several extensions arise naturally from the censoring framework
developed here.

\paragraph{Periodic inspection (Scheme~1).}
Under periodic inspection at times $\delta, 2\delta, 3\delta, \ldots$,
each component failure is localized to an inspection interval, yielding
component-level interval-censored data.  For a parallel system observed
under both Scheme~0 (system failure time~$T$) and Scheme~1 (component
states at inspection times), the likelihood for observation~$i$ takes
the form
\[
  L_i(\blam) = f_T(t_i;\blam) \prod_{j=1}^{m}
    \bigl[F_j(a_{ij}^+;\lambda_j) - F_j(a_{ij}^-;\lambda_j)\bigr],
\]
where $[a_{ij}^-,\, a_{ij}^+)$ is the inspection interval containing
$T_j$ (with the constraint $a_{ij}^+ \leq t_i$ for all~$j$ since
every component has failed by system failure time).  The interval
widths add information that resolves the left-censoring ambiguity of
Scheme~0, particularly for fast components.  Writing down and
evaluating this likelihood for exponential and Weibull components is a
natural next step; the Fisher information comparison across schemes
would quantify the value of periodic inspection.

\paragraph{Non-exponential components.}
The EM framework of \Cref{prop:em-algorithm} extends to any family where
the truncated conditional mean $\E[T_j \mid T_j < t]$ has a tractable
form.  For the Weibull family, the E-step requires the incomplete gamma
function, which is computationally available but not elementary.

\paragraph{Dependent components.}
Copula-based models could accommodate shared environmental stress or
manufacturing-induced dependence among components.  The censoring
equivalence (\Cref{prop:censoring-equiv}) continues to hold under
dependence---the system observation still comprises one exact and
$m-1$ left-censored component lifetimes---but the conditional
distributions in the E-step would no longer factor across components.
